\documentclass[11pt,a4paper]{article}

%====================================================
% PRÉAMBULE
%====================================================
\usepackage[utf8]{inputenc}
\usepackage[T1]{fontenc}
\usepackage[french]{babel}
\usepackage{lmodern}
\usepackage{geometry}
\geometry{margin=2cm}
\usepackage{amsmath,amssymb}
\usepackage{array,tabularx,multicol}
\usepackage{enumitem}
\usepackage{fancyhdr}
\usepackage{tikz}
\usepackage{tcolorbox}
\tcbuselibrary{skins,breakable}
\usepackage{hyperref}
\usepackage{siunitx}
\sisetup{locale=FR,detect-all}

\hypersetup{
    colorlinks=true,
    linkcolor=blue!60!black,
    urlcolor=blue!60!black
}

\pagestyle{fancy}
\fancyhf{}
\lhead{\textbf{TD complet -- Mathématiques 6e+}}
\rhead{\textbf{Nombres, grandeurs, géométrie, données}}
\cfoot{\thepage}

\setlength{\parindent}{0pt}
\setlength{\parskip}{0.45em}

%====================================================
% COULEURS ET BOÎTES
%====================================================
\definecolor{bleuprepa}{RGB}{25,70,140}
\definecolor{vertprepa}{RGB}{20,120,70}
\definecolor{rougeprepa}{RGB}{170,40,40}
\definecolor{orangeprepa}{RGB}{210,125,30}
\definecolor{violetprepa}{RGB}{100,60,150}
\definecolor{grisclair}{RGB}{245,247,250}

\tcbset{
  enhanced,
  breakable,
  sharp corners,
  boxrule=0.7pt,
  colback=white,
  fonttitle=\bfseries
}

\newtcolorbox{methodebox}{
  colframe=bleuprepa,
  colback=blue!3!white,
  title=Méthode
}

\newtcolorbox{retenirbox}{
  colframe=vertprepa,
  colback=green!3!white,
  title=À retenir
}

\newtcolorbox{erreurbox}{
  colframe=rougeprepa,
  colback=red!3!white,
  title=Erreur classique
}

\newtcolorbox{vigilancebox}{
  colframe=orangeprepa,
  colback=orange!5!white,
  title=Point de vigilance
}

\newtcolorbox{corrigebox}{
  colframe=violetprepa,
  colback=violet!3!white,
  title=Corrigé
}

\newcounter{exo}

\newenvironment{exercice}[3]{
\refstepcounter{exo}
\begin{tcolorbox}[
colframe=bleuprepa,
colback=grisclair,
title={Exercice \theexo{} -- #1 \hfill #2}
]
\textbf{Thème :} #3

\medskip
}{
\end{tcolorbox}
}

\newcommand{\diffA}{$\bigstar\ \star\ \star\ \star\ \star$}
\newcommand{\diffB}{$\bigstar\ \bigstar\ \star\ \star\ \star$}
\newcommand{\diffC}{$\bigstar\ \bigstar\ \bigstar\ \star\ \star$}
\newcommand{\diffD}{$\bigstar\ \bigstar\ \bigstar\ \bigstar\ \star$}
\newcommand{\diffE}{$\bigstar\ \bigstar\ \bigstar\ \bigstar\ \bigstar$}

\newcommand{\corrige}[2]{
\begin{corrigebox}
\textbf{Exercice #1.}

#2
\end{corrigebox}
}

%====================================================
% DOCUMENT
%====================================================
\begin{document}

%====================================================
% PAGE DE TITRE
%====================================================
\begin{titlepage}
\centering
\vspace*{1cm}

{\Huge \textbf{TD complet de mathématiques}}\\[0.4cm]
{\Large \textbf{Niveau Sixième +}}\\[0.6cm]

{\LARGE \textbf{Nombres, calculs, fractions, grandeurs, géométrie, données, probabilités et proportionnalité}}\\[1cm]

\begin{tcolorbox}[colframe=bleuprepa,colback=blue!3!white,width=0.9\textwidth,title=Objectif général]
Ce TD reprend les grandes méthodes du cours de sixième : numération décimale, opérations, fractions, pourcentages, grandeurs et mesures, géométrie plane, symétrie axiale, organisation de données, probabilités, proportionnalité et premiers raisonnements algébriques.

Il est conçu pour être progressif : les premiers exercices vérifient les bases, puis les situations deviennent plus longues, plus originales et plus exigeantes.
\end{tcolorbox}

\vfill

\begin{tabularx}{0.9\textwidth}{|>{\bfseries}l|X|}
\hline
Niveau & Sixième + \\
\hline
Durée indicative & Plusieurs séances : entraînement, approfondissement, devoir bilan \\
\hline
Compétences travaillées &
Calculer, raisonner, modéliser, représenter, chercher, communiquer, justifier une réponse. \\
\hline
Matériel utile & Règle graduée, compas, équerre, rapporteur, calculatrice ponctuellement, papier quadrillé. \\
\hline
\end{tabularx}

\vfill

{\large Document prêt à compiler sur Overleaf.}

\end{titlepage}

\tableofcontents
\newpage

%====================================================
\section{Rappels très courts de méthode}
%====================================================

\begin{retenirbox}
\textbf{Nombres décimaux.}
Dans un nombre décimal, chaque chiffre a une valeur qui dépend de son rang.
Par exemple :
\[
47,306 = 4\times 10 + 7\times 1 + 3\times 0,1 + 0\times 0,01 + 6\times 0,001.
\]
On peut comparer deux décimaux en alignant les virgules ou en ajoutant des zéros inutiles :
\[
3,7 = 3,70.
\]
\end{retenirbox}

\begin{methodebox}
\textbf{Comparer deux nombres décimaux.}
\begin{enumerate}[label=\arabic*)]
\item Comparer d'abord les parties entières.
\item Si elles sont égales, comparer les dixièmes.
\item Si les dixièmes sont égaux, comparer les centièmes.
\item Continuer si nécessaire.
\end{enumerate}
Exemple :
\[
12,405 < 12,45
\]
car
\[
12,405 = 12,405 \quad \text{et} \quad 12,45=12,450,
\]
or au rang des centièmes, $0<5$.
\end{methodebox}

\begin{retenirbox}
\textbf{Fractions.}
Pour deux entiers $a$ et $b$ avec $b\neq 0$,
\[
\frac{a}{b}
\]
désigne le quotient exact de $a$ par $b$ : c'est le nombre qui, multiplié par $b$, donne $a$.
Ainsi :
\[
b\times \frac{a}{b}=a.
\]
\end{retenirbox}

\begin{methodebox}
\textbf{Calculer une fraction d'une quantité.}
Pour calculer $\dfrac{a}{b}$ de $N$, on calcule :
\[
\dfrac{a}{b}\times N.
\]
Méthode pratique :
\begin{enumerate}[label=\arabic*)]
\item On divise $N$ par $b$.
\item On multiplie le résultat par $a$.
\end{enumerate}
Exemple :
\[
\frac{3}{4}\text{ de } 28 = 3\times (28\div 4)=3\times 7=21.
\]
\end{methodebox}

\begin{retenirbox}
\textbf{Pourcentages.}
\[
a\% = \frac{a}{100}.
\]
Donc $25\%$ signifie $\dfrac{25}{100}$, c'est-à-dire un quart.
\end{retenirbox}

\begin{methodebox}
\textbf{Proportionnalité.}
Deux grandeurs sont proportionnelles lorsqu'on passe toujours de l'une à l'autre en multipliant par un même nombre non nul.

Exemple : si $3$ kg de pommes coûtent $7,50$ euros, alors $1$ kg coûte :
\[
7,50\div 3 = 2,50 \text{ euros}.
\]
C'est le retour à l'unité.
\end{methodebox}

\begin{retenirbox}
\textbf{Périmètres et aires.}
\[
P_{\text{rectangle}}=2\times(L+\ell),
\qquad
A_{\text{rectangle}}=L\times \ell.
\]
\[
P_{\text{cercle}}=\pi\times D=2\times \pi\times R.
\]
\[
A_{\text{carré}}=c\times c.
\]
\end{retenirbox}

\begin{retenirbox}
\textbf{Angles et triangles.}
La somme des mesures des trois angles d'un triangle vaut :
\[
180^\circ.
\]
Dans un triangle équilatéral, les trois angles mesurent :
\[
60^\circ.
\]
\end{retenirbox}

\begin{retenirbox}
\textbf{Probabilités simples.}
Dans une situation d'équiprobabilité :
\[
P(\text{événement})=\frac{\text{nombre de cas favorables}}{\text{nombre de cas possibles}}.
\]
Une probabilité est toujours comprise entre $0$ et $1$.
\end{retenirbox}

\begin{erreurbox}
\begin{itemize}
\item Multiplier ne veut pas toujours dire rendre plus grand : multiplier par $0,1$ revient à diviser par $10$.
\item $3,9$ n'est pas plus petit que $3,12$ : on compare $3,90$ et $3,12$.
\item Dans une aire, on utilise des unités carrées : $\text{cm}^2$, $\text{m}^2$, etc.
\item Dans une probabilité, on ne met pas seulement le nombre de cas favorables : il faut le comparer au nombre de cas possibles.
\end{itemize}
\end{erreurbox}

\newpage

%====================================================
\section{Exercices d'application directe}
%====================================================

\begin{exercice}{La fusée des décimaux}{\diffA}{Numération décimale}
Une fusée d'exploration affiche l'altitude suivante :
\[
12\,407,536 \text{ km}.
\]
\begin{enumerate}[label=\arabic*)]
\item Donner le chiffre des unités.
\item Donner le chiffre des dixièmes.
\item Donner le chiffre des millièmes.
\item Donner le nombre de kilomètres entiers.
\item Écrire ce nombre sous la forme :
\[
12\,000+400+\cdots
\]
\end{enumerate}
\end{exercice}

\begin{exercice}{Classement de dragons}{\diffA}{Comparer des nombres décimaux}
Quatre dragons mesurent :
\[
7,8\text{ m},\quad 7,08\text{ m},\quad 7,805\text{ m},\quad 7,85\text{ m}.
\]
\begin{enumerate}[label=\arabic*)]
\item Ranger ces longueurs dans l'ordre croissant.
\item Quel dragon est le plus grand ?
\item Quel dragon est le plus petit ?
\end{enumerate}
\end{exercice}

\begin{exercice}{Tickets de fête foraine}{\diffA}{Addition et soustraction de décimaux}
Lina achète :
\begin{itemize}
\item une barbe à papa à $2,40$ euros ;
\item un tour de grande roue à $3,75$ euros ;
\item un jeu d'adresse à $1,90$ euro.
\end{itemize}
Elle avait $10$ euros.
\begin{enumerate}[label=\arabic*)]
\item Combien dépense-t-elle ?
\item Combien lui reste-t-il ?
\item Vérifier le résultat par un ordre de grandeur.
\end{enumerate}
\end{exercice}

\begin{exercice}{Fractions-pizza}{\diffA}{Fractions simples}
Une pizza est découpée en $8$ parts égales.
\begin{enumerate}[label=\arabic*)]
\item Quelle fraction de la pizza représente une part ?
\item Tom mange $3$ parts. Quelle fraction de la pizza a-t-il mangée ?
\item Quelle fraction reste-t-il ?
\item Écrire la fraction mangée sous forme décimale.
\end{enumerate}
\end{exercice}

\begin{exercice}{Le goûter des robots}{\diffA}{Fraction d'une quantité}
Un robot doit distribuer $36$ biscuits.
Il en donne $\dfrac{2}{3}$ à une classe.
\begin{enumerate}[label=\arabic*)]
\item Combien de biscuits distribue-t-il ?
\item Combien lui reste-t-il ?
\end{enumerate}
\end{exercice}

\begin{exercice}{Conversions express}{\diffA}{Grandeurs et mesures}
Convertir :
\begin{enumerate}[label=\arabic*)]
\item $3$ m en cm.
\item $450$ cm en m.
\item $2,5$ L en cL.
\item $75$ min en h et min.
\item $3,2$ m$^2$ en dm$^2$.
\end{enumerate}
\end{exercice}

\begin{exercice}{Le rectangle invisible}{\diffA}{Périmètre et aire}
Un rectangle a une longueur de $8$ cm et une largeur de $5$ cm.
\begin{enumerate}[label=\arabic*)]
\item Calculer son périmètre.
\item Calculer son aire.
\item Préciser les unités utilisées.
\end{enumerate}
\end{exercice}

\begin{exercice}{Dé truqué ou pas ?}{\diffA}{Probabilités simples}
On lance un dé équilibré à six faces numérotées de $1$ à $6$.
\begin{enumerate}[label=\arabic*)]
\item Quelle est la probabilité d'obtenir $4$ ?
\item Quelle est la probabilité d'obtenir un nombre pair ?
\item Quelle est la probabilité d'obtenir un nombre supérieur à $6$ ?
\item Quelle est la probabilité d'obtenir un nombre inférieur ou égal à $6$ ?
\end{enumerate}
\end{exercice}

\newpage

%====================================================
\section{Exercices de méthode et exercices standards}
%====================================================

\begin{exercice}{Les pièces du trésor}{\diffB}{Multiplication de décimaux}
Un coffre contient $12$ sacs.
Chaque sac contient $3,5$ kg de pièces.
\begin{enumerate}[label=\arabic*)]
\item Calculer la masse totale des pièces.
\item Expliquer pourquoi le résultat est cohérent avec l'ordre de grandeur $12\times 4$.
\item Si chaque kilogramme de pièces vaut $18,60$ euros, calculer la valeur totale du trésor.
\end{enumerate}
\end{exercice}

\begin{exercice}{Potion magique dosée}{\diffB}{Division décimale}
Une potion de $1,5$ L doit être versée équitablement dans $6$ fioles.
\begin{enumerate}[label=\arabic*)]
\item Convertir $1,5$ L en cL.
\item Combien de cL contient chaque fiole ?
\item Combien de mL contient chaque fiole ?
\end{enumerate}
\end{exercice}

\begin{exercice}{Bus pour une sortie scolaire}{\diffB}{Division euclidienne}
Un collège organise une sortie avec $286$ élèves.
Un bus peut transporter $48$ élèves.
\begin{enumerate}[label=\arabic*)]
\item Effectuer la division euclidienne de $286$ par $48$.
\item Combien de bus totalement remplis peut-on avoir ?
\item Combien de bus faut-il prévoir au minimum ?
\item Combien de places seront vides dans le dernier bus si tous les élèves partent ?
\end{enumerate}
\end{exercice}

\begin{exercice}{Le laboratoire des fractions}{\diffB}{Égalités de fractions}
\begin{enumerate}[label=\arabic*)]
\item Compléter :
\[
\frac{2}{3}=\frac{\cdots}{12}.
\]
\item Compléter :
\[
\frac{5}{7}=\frac{25}{\cdots}.
\]
\item Dire si les fractions suivantes sont égales :
\[
\frac{6}{8}\quad \text{et}\quad \frac{9}{12}.
\]
\item Simplifier :
\[
\frac{18}{24}.
\]
\end{enumerate}
\end{exercice}

\begin{exercice}{Le gâteau des anniversaires}{\diffB}{Addition de fractions}
Maya mange $\dfrac{1}{4}$ d'un gâteau.
Noé mange $\dfrac{1}{3}$ du même gâteau.
\begin{enumerate}[label=\arabic*)]
\item Quelle fraction du gâteau ont-ils mangée ensemble ?
\item Quelle fraction reste-t-il ?
\item Qui a mangé la plus grande part ? Justifier.
\end{enumerate}
\end{exercice}

\begin{exercice}{Soldes au magasin de potions}{\diffB}{Pourcentages}
Une cape coûte $80$ euros.
Elle est soldée à $25\%$.
\begin{enumerate}[label=\arabic*)]
\item Calculer le montant de la réduction.
\item Calculer le nouveau prix.
\item Expliquer pourquoi $25\%$ correspond à un quart.
\end{enumerate}
\end{exercice}

\begin{exercice}{Le cercle du magicien}{\diffB}{Cercle et périmètre}
Un magicien trace un cercle de rayon $4$ cm.
On prendra $\pi\approx 3,14$.
\begin{enumerate}[label=\arabic*)]
\item Quel est le diamètre du cercle ?
\item Calculer le périmètre du cercle.
\item Donner un arrondi au dixième de centimètre.
\end{enumerate}
\end{exercice}

\begin{exercice}{Angles dans un triangle}{\diffB}{Somme des angles d'un triangle}
Dans un triangle $ABC$, on sait que :
\[
\widehat{A}=42^\circ
\quad\text{et}\quad
\widehat{B}=68^\circ.
\]
\begin{enumerate}[label=\arabic*)]
\item Calculer la mesure de l'angle $\widehat{C}$.
\item Le triangle peut-il être rectangle ? Justifier.
\end{enumerate}
\end{exercice}

\newpage

%====================================================
\section{Exercices intermédiaires}
%====================================================

\begin{exercice}{Le parc des pandas roux}{\diffC}{Tableaux, proportionnalité et pourcentages}
Dans un parc animalier, on observe $40$ pandas roux.
Parmi eux, $18$ sont des femelles, $22$ sont des mâles, et $10$ ont moins d'un an.
\begin{enumerate}[label=\arabic*)]
\item Quelle fraction des pandas sont des femelles ?
\item Exprimer cette proportion en pourcentage.
\item Quelle fraction des pandas ont moins d'un an ?
\item Le soigneur dit : « Un quart des pandas ont moins d'un an. » A-t-il raison ?
\end{enumerate}
\end{exercice}

\begin{exercice}{Le plan secret du château}{\diffC}{Échelles}
Sur le plan d'un château, $1$ cm représente $5$ m dans la réalité.
\begin{enumerate}[label=\arabic*)]
\item Une salle mesure $3,4$ cm sur le plan. Quelle est sa longueur réelle ?
\item Un couloir mesure $27,5$ m dans la réalité. Quelle longueur mesure-t-il sur le plan ?
\item Une cour rectangulaire mesure $4$ cm par $2,5$ cm sur le plan. Calculer ses dimensions réelles.
\item Calculer l'aire réelle de cette cour.
\end{enumerate}
\end{exercice}

\begin{exercice}{Le robot peintre}{\diffC}{Aire, périmètre et coût}
Un robot peint une fresque rectangulaire de longueur $6,5$ m et de largeur $2,4$ m.
Un pot de peinture couvre $4$ m$^2$.
Un pot coûte $13,80$ euros.
\begin{enumerate}[label=\arabic*)]
\item Calculer l'aire de la fresque.
\item Combien de pots sont nécessaires ?
\item Calculer le prix total.
\item Pourquoi ne peut-on pas acheter $3,9$ pots ?
\end{enumerate}
\end{exercice}

\begin{exercice}{Le triangle mystérieux}{\diffC}{Triangles particuliers}
Un triangle $ABC$ est isocèle en $A$.
On sait que :
\[
\widehat{A}=40^\circ.
\]
\begin{enumerate}[label=\arabic*)]
\item Que peut-on dire des angles $\widehat{B}$ et $\widehat{C}$ ?
\item Calculer $\widehat{B}$ et $\widehat{C}$.
\item Ce triangle est-il équilatéral ? Justifier.
\end{enumerate}
\end{exercice}

\begin{exercice}{La médiatrice du trésor}{\diffC}{Médiatrice}
Deux pirates enterrent un trésor à égale distance de deux palmiers $A$ et $B$.
La distance entre les palmiers est de $12$ m.
\begin{enumerate}[label=\arabic*)]
\item Sur quelle ligne doit se trouver le trésor ? Donner le nom mathématique.
\item Expliquer pourquoi tous les points de cette ligne sont à égale distance de $A$ et $B$.
\item Si le trésor est à $8$ m de $A$, quelle est sa distance à $B$ ?
\end{enumerate}
\end{exercice}

\begin{exercice}{Le sac de billes de l'alchimiste}{\diffC}{Probabilités}
Un sac contient $5$ billes rouges, $3$ billes bleues et $2$ billes vertes.
On tire une bille au hasard.
\begin{enumerate}[label=\arabic*)]
\item Combien y a-t-il de billes au total ?
\item Quelle est la probabilité de tirer une bille rouge ?
\item Quelle est la probabilité de tirer une bille bleue ou verte ?
\item Exprimer cette dernière probabilité en pourcentage.
\end{enumerate}
\end{exercice}

\begin{exercice}{Programme de calcul du petit monstre}{\diffC}{Pensée algébrique}
Un petit monstre choisit un nombre.
Il applique le programme suivant :
\begin{enumerate}[label=\arabic*)]
\item Multiplier par $3$.
\item Ajouter $7$.
\item Doubler le résultat.
\end{enumerate}
\begin{enumerate}[label=\alph*)]
\item Quel résultat obtient-il s'il choisit $4$ ?
\item Quel résultat obtient-il s'il choisit $10$ ?
\item Il obtient $50$. Quel nombre avait-il choisi au départ ?
\end{enumerate}
\end{exercice}

\begin{exercice}{La course des escargots turbo}{\diffC}{Proportionnalité}
Un escargot turbo parcourt $18$ cm en $6$ minutes à vitesse constante.
\begin{enumerate}[label=\arabic*)]
\item Quelle distance parcourt-il en $1$ minute ?
\item Quelle distance parcourt-il en $15$ minutes ?
\item Combien de temps met-il pour parcourir $45$ cm ?
\item Expliquer pourquoi cette situation est proportionnelle.
\end{enumerate}
\end{exercice}

\newpage

%====================================================
\section{Exercices difficiles}
%====================================================

\begin{exercice}{Le banquet des fractions}{\diffD}{Fractions, proportionnalité et raisonnement}
Pour un banquet, on prépare $120$ tartelettes.
\[
\frac{1}{3}
\]
sont aux pommes,
\[
\frac{2}{5}
\]
sont au chocolat, et le reste est à la fraise.
\begin{enumerate}[label=\arabic*)]
\item Calculer le nombre de tartelettes aux pommes.
\item Calculer le nombre de tartelettes au chocolat.
\item Calculer le nombre de tartelettes à la fraise.
\item Quelle fraction du total représentent les tartelettes à la fraise ?
\item Vérifier que la somme des trois fractions vaut $1$.
\end{enumerate}
\end{exercice}

\begin{exercice}{Le mini-golf circulaire}{\diffD}{Périmètres composés}
Un parcours de mini-golf est formé d'un rectangle de longueur $12$ m et de largeur $5$ m, prolongé à chaque extrémité par un demi-cercle de diamètre $5$ m.
On prendra $\pi\approx 3,14$.
\begin{enumerate}[label=\arabic*)]
\item Faire un schéma à main levée.
\item Expliquer pourquoi les deux demi-cercles forment ensemble un cercle complet.
\item Calculer la longueur totale du contour.
\item Donner un arrondi au centimètre.
\end{enumerate}
\end{exercice}

\begin{exercice}{Le coffre codé}{\diffD}{Décimaux, fractions et logique}
Pour ouvrir un coffre, il faut trouver un code à trois nombres.
\begin{itemize}
\item Le premier nombre est le résultat de $4,8\times 0,5$.
\item Le deuxième nombre est égal à $\dfrac{3}{4}$ de $60$.
\item Le troisième nombre est le plus petit nombre parmi :
\[
7,09;\quad 7,9;\quad 7,099;\quad 7,901.
\]
\end{itemize}
\begin{enumerate}[label=\arabic*)]
\item Trouver les trois nombres du code.
\item Justifier chaque réponse.
\item Ranger les quatre nombres proposés pour le troisième code dans l'ordre croissant.
\end{enumerate}
\end{exercice}

\begin{exercice}{Une enquête au collège}{\diffD}{Organisation de données et pourcentages}
Dans un collège de $600$ élèves, une enquête est menée sur le moyen de transport principal.
\[
\begin{array}{|c|c|}
\hline
\text{Moyen de transport} & \text{Nombre d'élèves}\\
\hline
\text{Bus} & 240\\
\hline
\text{Voiture} & 150\\
\hline
\text{Vélo} & 90\\
\hline
\text{Marche} & 120\\
\hline
\end{array}
\]
\begin{enumerate}[label=\arabic*)]
\item Vérifier que le total est cohérent.
\item Calculer le pourcentage d'élèves venant en bus.
\item Calculer le pourcentage d'élèves venant à vélo.
\item Quel moyen de transport représente exactement un quart des élèves ?
\item Présenter les résultats dans un tableau avec une colonne « pourcentage ».
\end{enumerate}
\end{exercice}

\begin{exercice}{Le jardin symétrique}{\diffD}{Symétrie axiale et aire}
Un jardin rectangulaire mesure $10$ m sur $6$ m.
Une allée droite partage le jardin en deux parties symétriques selon l'axe de symétrie passant par les milieux des deux grands côtés.
\begin{enumerate}[label=\arabic*)]
\item Faire un schéma.
\item Calculer l'aire totale du jardin.
\item Quelle est l'aire de chaque moitié ?
\item Expliquer pourquoi les deux parties ont la même aire.
\item Si on plante $4$ fleurs par m$^2$ dans une moitié, combien faut-il de fleurs ?
\end{enumerate}
\end{exercice}

\newpage

%====================================================
\section{Exercices ambitieux, originaux et de synthèse}
%====================================================

\begin{exercice}{Mission sur Mars}{\diffE}{Synthèse : décimaux, proportionnalité, durées}
Un petit robot martien avance à vitesse constante.
Il parcourt $7,5$ m en $3$ minutes.
Sa batterie lui permet de fonctionner pendant $48$ minutes.
\begin{enumerate}[label=\arabic*)]
\item Quelle distance parcourt-il en $1$ minute ?
\item Quelle distance peut-il parcourir avec une batterie pleine ?
\item Le robot doit rejoindre une roche située à $125$ m. Peut-il l'atteindre avec une seule batterie ?
\item Combien de minutes lui manque-t-il ou lui reste-t-il ?
\item Proposer une phrase de conclusion comme dans une copie de contrôle.
\end{enumerate}
\end{exercice}

\begin{exercice}{Le festival des formes}{\diffE}{Synthèse : géométrie, périmètre, aire, coût}
Une scène de festival a la forme d'un rectangle de longueur $14$ m et de largeur $8$ m.
On veut poser une barrière tout autour, sauf sur une ouverture de $3$ m.
On veut aussi recouvrir le sol avec des dalles.
Chaque dalle couvre $2$ m$^2$ et coûte $9,50$ euros.
Un mètre de barrière coûte $6,20$ euros.
\begin{enumerate}[label=\arabic*)]
\item Calculer le périmètre complet de la scène.
\item Calculer la longueur de barrière nécessaire.
\item Calculer le prix de la barrière.
\item Calculer l'aire de la scène.
\item Calculer le nombre de dalles nécessaires.
\item Calculer le prix total des dalles.
\item Calculer le budget total.
\end{enumerate}
\end{exercice}

\begin{exercice}{Le grand laboratoire des probabilités}{\diffE}{Synthèse : probabilités, fractions, pourcentages}
Un laboratoire teste des cristaux magiques.
Une boîte contient :
\[
8 \text{ cristaux bleus},\quad 6 \text{ cristaux rouges},\quad 4 \text{ cristaux verts},\quad 2 \text{ cristaux dorés}.
\]
On tire un cristal au hasard.
\begin{enumerate}[label=\arabic*)]
\item Calculer le nombre total de cristaux.
\item Calculer la probabilité de tirer un cristal doré.
\item Calculer la probabilité de tirer un cristal bleu.
\item Calculer la probabilité de ne pas tirer un cristal vert.
\item Écrire chaque probabilité sous forme de fraction simplifiée, de nombre décimal et de pourcentage.
\item Classer les événements du moins probable au plus probable :
\[
D=\text{« tirer doré »},\quad
B=\text{« tirer bleu »},\quad
N=\text{« ne pas tirer vert »}.
\]
\end{enumerate}
\end{exercice}

%====================================================
\section{Deux grands exercices bilan}
%====================================================

\begin{exercice}{Bilan 1 -- La sortie au parc aquatique}{\diffD}{Synthèse complète}
Une classe de $28$ élèves organise une sortie au parc aquatique.
Le billet d'entrée coûte $12,50$ euros par élève.
Le transport en bus coûte $180$ euros pour toute la classe.
Le professeur achète aussi $28$ goûters à $2,40$ euros chacun.

Le parc possède un bassin rectangulaire de longueur $25$ m et de largeur $12$ m.
Une zone circulaire de jeux a un diamètre de $8$ m.
On prendra $\pi\approx 3,14$.
\begin{enumerate}[label=\arabic*)]
\item Calculer le prix total des billets.
\item Calculer le prix total des goûters.
\item Calculer le coût total de la sortie.
\item Calculer le coût moyen par élève.
\item Calculer l'aire du bassin rectangulaire.
\item Calculer le périmètre de la zone circulaire.
\item Le responsable dit que la zone circulaire a un contour d'environ $25$ m. A-t-il raison ?
\end{enumerate}
\end{exercice}

\begin{exercice}{Bilan 2 -- Le village médiéval}{\diffE}{Synthèse complète et raisonnement}
Un village médiéval organise une fête.
Il y a $360$ visiteurs.
\[
\frac{2}{5}
\]
des visiteurs participent au tir à l'arc.
$30\%$ des visiteurs visitent le château.
Les autres vont au marché.

Le plan du marché est à l'échelle suivante :
\[
1\text{ cm sur le plan représente }4\text{ m dans la réalité}.
\]
Sur le plan, la place du marché est un rectangle de $6$ cm sur $3,5$ cm.
\begin{enumerate}[label=\arabic*)]
\item Calculer le nombre de visiteurs au tir à l'arc.
\item Calculer le nombre de visiteurs au château.
\item Calculer le nombre de visiteurs au marché.
\item Calculer les dimensions réelles de la place du marché.
\item Calculer l'aire réelle de la place du marché.
\item On installe un stand pour $12$ m$^2$. Combien de stands peut-on installer au maximum ?
\item Proposer une conclusion expliquant si la place est suffisante pour installer $25$ stands.
\end{enumerate}
\end{exercice}

\newpage

%====================================================
\section{Corrigé complet et détaillé}
%====================================================

\corrige{1}{
\begin{enumerate}[label=\arabic*)]
\item Dans $12\,407,536$, le chiffre des unités est $7$.
\item Le chiffre des dixièmes est le premier chiffre après la virgule : $5$.
\item Le chiffre des millièmes est le troisième chiffre après la virgule : $6$.
\item Le nombre de kilomètres entiers est $12\,407$.
\item On décompose :
\[
12\,407,536=10\,000+2\,000+400+7+0,5+0,03+0,006.
\]
Chaque terme correspond à la valeur d'un chiffre selon son rang.
\end{enumerate}
}

\corrige{2}{
On compare les nombres en ajoutant des zéros :
\[
7,8=7,800,\quad 7,08=7,080,\quad 7,805=7,805,\quad 7,85=7,850.
\]
Donc :
\[
7,080<7,800<7,805<7,850.
\]
L'ordre croissant est :
\[
7,08<7,8<7,805<7,85.
\]
Le plus grand dragon mesure $7,85$ m.
Le plus petit mesure $7,08$ m.
}

\corrige{3}{
\begin{enumerate}[label=\arabic*)]
\item La dépense totale est :
\[
2,40+3,75+1,90=8,05.
\]
Lina dépense donc $8,05$ euros.
\item Il lui reste :
\[
10-8,05=1,95.
\]
Il lui reste $1,95$ euro.
\item Ordre de grandeur :
\[
2,40\approx 2,\quad 3,75\approx 4,\quad 1,90\approx 2.
\]
Donc la dépense est proche de $2+4+2=8$ euros, ce qui est cohérent avec $8,05$ euros.
\end{enumerate}
}

\corrige{4}{
\begin{enumerate}[label=\arabic*)]
\item Une part représente :
\[
\frac{1}{8}.
\]
\item Tom mange $3$ parts, donc :
\[
\frac{3}{8}.
\]
\item Il reste :
\[
\frac{8}{8}-\frac{3}{8}=\frac{5}{8}.
\]
\item Sous forme décimale :
\[
\frac{3}{8}=3\div 8=0,375.
\]
Tom a mangé $0,375$ pizza, c'est-à-dire $37,5\%$ de la pizza.
\end{enumerate}
}

\corrige{5}{
On cherche $\dfrac{2}{3}$ de $36$.
\[
\frac{2}{3}\times 36=2\times(36\div 3)=2\times 12=24.
\]
Le robot distribue $24$ biscuits.

Il lui reste :
\[
36-24=12.
\]
Il reste donc $12$ biscuits.
}

\corrige{6}{
\begin{enumerate}[label=\arabic*)]
\item $1$ m $=100$ cm, donc :
\[
3\text{ m}=300\text{ cm}.
\]
\item $450$ cm $=4,5$ m.
\item $1$ L $=100$ cL, donc :
\[
2,5\text{ L}=250\text{ cL}.
\]
\item $75$ min $=60$ min $+15$ min, donc :
\[
75\text{ min}=1\text{ h }15\text{ min}.
\]
\item $1$ m$^2=100$ dm$^2$, donc :
\[
3,2\text{ m}^2=320\text{ dm}^2.
\]
\end{enumerate}
}

\corrige{7}{
\begin{enumerate}[label=\arabic*)]
\item Le périmètre d'un rectangle est :
\[
P=2\times(L+\ell).
\]
Donc :
\[
P=2\times(8+5)=2\times 13=26.
\]
Le périmètre est $26$ cm.
\item L'aire d'un rectangle est :
\[
A=L\times \ell.
\]
Donc :
\[
A=8\times 5=40.
\]
L'aire est $40$ cm$^2$.
\item Le périmètre s'exprime en cm, l'aire en cm$^2$.
\end{enumerate}
}

\corrige{8}{
Le dé possède $6$ issues équiprobables.
\begin{enumerate}[label=\arabic*)]
\item Une seule face porte le numéro $4$, donc :
\[
P(4)=\frac{1}{6}.
\]
\item Les nombres pairs sont $2$, $4$, $6$, donc $3$ cas favorables :
\[
P(\text{pair})=\frac{3}{6}=\frac{1}{2}.
\]
\item Il est impossible d'obtenir un nombre supérieur à $6$ :
\[
P=0.
\]
\item Il est certain d'obtenir un nombre inférieur ou égal à $6$ :
\[
P=1.
\]
\end{enumerate}
}

\corrige{9}{
\begin{enumerate}[label=\arabic*)]
\item La masse totale est :
\[
12\times 3,5=42.
\]
Le coffre contient $42$ kg de pièces.
\item Comme $3,5$ est proche de $4$, on peut estimer :
\[
12\times 4=48.
\]
Le résultat $42$ kg est proche de $48$ kg : il est cohérent.
\item La valeur totale est :
\[
42\times 18,60=781,20.
\]
Le trésor vaut donc $781,20$ euros.
\end{enumerate}
}

\corrige{10}{
\begin{enumerate}[label=\arabic*)]
\item $1$ L $=100$ cL, donc :
\[
1,5\text{ L}=150\text{ cL}.
\]
\item On partage $150$ cL en $6$ fioles :
\[
150\div 6=25.
\]
Chaque fiole contient $25$ cL.
\item $1$ cL $=10$ mL, donc :
\[
25\text{ cL}=250\text{ mL}.
\]
Chaque fiole contient $250$ mL.
\end{enumerate}
}

\corrige{11}{
\begin{enumerate}[label=\arabic*)]
\item Division euclidienne :
\[
286=48\times 5+46.
\]
\item On peut remplir totalement $5$ bus.
\item Il reste $46$ élèves, donc il faut un bus supplémentaire.
Il faut donc $6$ bus au minimum.
\item Avec $6$ bus, il y a :
\[
6\times 48=288
\]
places.
Il y aura :
\[
288-286=2
\]
places vides.
\end{enumerate}
}

\corrige{12}{
\begin{enumerate}[label=\arabic*)]
\item Pour passer de $3$ à $12$, on multiplie par $4$, donc :
\[
\frac{2}{3}=\frac{8}{12}.
\]
\item Pour passer de $5$ à $25$, on multiplie par $5$, donc on multiplie aussi le dénominateur par $5$ :
\[
\frac{5}{7}=\frac{25}{35}.
\]
\item On simplifie :
\[
\frac{6}{8}=\frac{3}{4},
\qquad
\frac{9}{12}=\frac{3}{4}.
\]
Les deux fractions sont donc égales.
\item
\[
\frac{18}{24}=\frac{3}{4}
\]
car on divise le numérateur et le dénominateur par $6$.
\end{enumerate}
}

\corrige{13}{
\begin{enumerate}[label=\arabic*)]
\item On additionne :
\[
\frac{1}{4}+\frac{1}{3}.
\]
Un dénominateur commun est $12$ :
\[
\frac{1}{4}=\frac{3}{12},
\qquad
\frac{1}{3}=\frac{4}{12}.
\]
Donc :
\[
\frac{1}{4}+\frac{1}{3}=\frac{7}{12}.
\]
Ils ont mangé $\dfrac{7}{12}$ du gâteau.
\item Il reste :
\[
\frac{12}{12}-\frac{7}{12}=\frac{5}{12}.
\]
\item On compare :
\[
\frac{1}{4}=\frac{3}{12},
\qquad
\frac{1}{3}=\frac{4}{12}.
\]
Comme $4>3$, Noé a mangé la plus grande part.
\end{enumerate}
}

\corrige{14}{
\begin{enumerate}[label=\arabic*)]
\item $25\%=\dfrac{25}{100}=\dfrac{1}{4}$.
La réduction vaut donc :
\[
80\div 4=20.
\]
La réduction est de $20$ euros.
\item Le nouveau prix est :
\[
80-20=60.
\]
La cape coûte $60$ euros.
\item $25\%=\dfrac{25}{100}$. En divisant par $25$ :
\[
\frac{25}{100}=\frac{1}{4}.
\]
Donc $25\%$ correspond à un quart.
\end{enumerate}
}

\corrige{15}{
\begin{enumerate}[label=\arabic*)]
\item Le rayon est $4$ cm, donc le diamètre est :
\[
D=2\times 4=8.
\]
Le diamètre est $8$ cm.
\item Le périmètre d'un cercle est :
\[
P=\pi\times D.
\]
Donc :
\[
P\approx 3,14\times 8=25,12.
\]
\item Arrondi au dixième :
\[
25,12\approx 25,1.
\]
Le périmètre est environ $25,1$ cm.
\end{enumerate}
}

\corrige{16}{
\begin{enumerate}[label=\arabic*)]
\item Dans un triangle, la somme des angles vaut $180^\circ$.
Donc :
\[
\widehat{C}=180^\circ-42^\circ-68^\circ=70^\circ.
\]
\item Aucun angle ne mesure $90^\circ$.
Le triangle n'est donc pas rectangle.
\end{enumerate}
}

\corrige{17}{
\begin{enumerate}[label=\arabic*)]
\item La fraction de femelles est :
\[
\frac{18}{40}.
\]
On peut simplifier :
\[
\frac{18}{40}=\frac{9}{20}.
\]
\item Pour l'exprimer en pourcentage :
\[
\frac{9}{20}=\frac{45}{100}=45\%.
\]
\item La fraction de pandas ayant moins d'un an est :
\[
\frac{10}{40}=\frac{1}{4}.
\]
\item Oui, le soigneur a raison car :
\[
\frac{1}{4}=25\%.
\]
\end{enumerate}
}

\corrige{18}{
\begin{enumerate}[label=\arabic*)]
\item $1$ cm représente $5$ m.
Donc :
\[
3,4\times 5=17.
\]
La salle mesure $17$ m.
\item On cherche combien de fois $5$ m contient $27,5$ m :
\[
27,5\div 5=5,5.
\]
Le couloir mesure $5,5$ cm sur le plan.
\item Dimensions réelles :
\[
4\times 5=20\text{ m},
\qquad
2,5\times 5=12,5\text{ m}.
\]
\item Aire réelle :
\[
A=20\times 12,5=250.
\]
L'aire réelle de la cour est $250$ m$^2$.
\end{enumerate}
}

\corrige{19}{
\begin{enumerate}[label=\arabic*)]
\item Aire :
\[
A=6,5\times 2,4=15,6.
\]
La fresque a une aire de $15,6$ m$^2$.
\item Un pot couvre $4$ m$^2$ :
\[
15,6\div 4=3,9.
\]
Il faut donc acheter $4$ pots.
\item Prix :
\[
4\times 13,80=55,20.
\]
Le prix total est $55,20$ euros.
\item On ne peut pas acheter $3,9$ pots car les pots sont vendus entiers. Il faut arrondir au nombre entier supérieur.
\end{enumerate}
}

\corrige{20}{
\begin{enumerate}[label=\arabic*)]
\item Dans un triangle isocèle en $A$, les côtés $AB$ et $AC$ sont égaux, donc les angles à la base sont égaux :
\[
\widehat{B}=\widehat{C}.
\]
\item La somme des angles vaut $180^\circ$ :
\[
\widehat{B}+\widehat{C}=180^\circ-40^\circ=140^\circ.
\]
Comme ils sont égaux :
\[
\widehat{B}=\widehat{C}=70^\circ.
\]
\item Un triangle équilatéral a trois angles de $60^\circ$.
Ici les angles sont $40^\circ$, $70^\circ$, $70^\circ$.
Il n'est donc pas équilatéral.
\end{enumerate}
}

\corrige{21}{
\begin{enumerate}[label=\arabic*)]
\item Le trésor doit se trouver sur la médiatrice du segment $[AB]$.
\item La médiatrice d'un segment est l'ensemble des points équidistants des extrémités de ce segment.
Donc tout point de la médiatrice de $[AB]$ est à la même distance de $A$ et de $B$.
\item Si le trésor est à $8$ m de $A$, alors il est aussi à $8$ m de $B$.
\end{enumerate}
}

\corrige{22}{
\begin{enumerate}[label=\arabic*)]
\item Total :
\[
5+3+2=10.
\]
Il y a $10$ billes.
\item Probabilité de rouge :
\[
P(R)=\frac{5}{10}=\frac{1}{2}.
\]
\item Billes bleues ou vertes :
\[
3+2=5.
\]
Donc :
\[
P(B\text{ ou }V)=\frac{5}{10}=\frac{1}{2}.
\]
\item
\[
\frac{1}{2}=0,5=50\%.
\]
\end{enumerate}
}

\corrige{23}{
\begin{enumerate}[label=\alph*)]
\item Avec $4$ :
\[
4\times 3=12,
\quad
12+7=19,
\quad
19\times 2=38.
\]
Le résultat est $38$.
\item Avec $10$ :
\[
10\times 3=30,
\quad
30+7=37,
\quad
37\times 2=74.
\]
Le résultat est $74$.
\item On remonte le programme à l'envers.
Le dernier calcul était « doubler ». Donc avant de doubler, on avait :
\[
50\div 2=25.
\]
Avant d'ajouter $7$, on avait :
\[
25-7=18.
\]
Avant de multiplier par $3$, on avait :
\[
18\div 3=6.
\]
Le nombre choisi était $6$.
\end{enumerate}
}

\corrige{24}{
\begin{enumerate}[label=\arabic*)]
\item En $6$ minutes, il parcourt $18$ cm.
En $1$ minute :
\[
18\div 6=3.
\]
Il parcourt $3$ cm par minute.
\item En $15$ minutes :
\[
15\times 3=45.
\]
Il parcourt $45$ cm.
\item Pour parcourir $45$ cm :
\[
45\div 3=15.
\]
Il met $15$ minutes.
\item La situation est proportionnelle car la distance parcourue est toujours obtenue en multipliant la durée par $3$.
\end{enumerate}
}

\corrige{25}{
\begin{enumerate}[label=\arabic*)]
\item Tartelettes aux pommes :
\[
\frac{1}{3}\times 120=120\div 3=40.
\]
\item Tartelettes au chocolat :
\[
\frac{2}{5}\times 120=2\times(120\div 5)=2\times 24=48.
\]
\item Tartelettes à la fraise :
\[
120-40-48=32.
\]
\item Fraction à la fraise :
\[
\frac{32}{120}=\frac{4}{15}.
\]
\item Vérifions :
\[
\frac{1}{3}+\frac{2}{5}+\frac{4}{15}
=
\frac{5}{15}+\frac{6}{15}+\frac{4}{15}
=
\frac{15}{15}=1.
\]
La répartition est cohérente.
\end{enumerate}
}

\corrige{26}{
\begin{enumerate}[label=\arabic*)]
\item Le schéma est un rectangle de longueur $12$ m avec deux demi-cercles aux extrémités.
\item Deux demi-cercles de même diamètre forment un cercle complet.
\item Le contour est composé de deux longueurs droites de $12$ m et du périmètre d'un cercle de diamètre $5$ m :
\[
P=12+12+\pi\times 5.
\]
Donc :
\[
P\approx 24+3,14\times 5=24+15,7=39,7.
\]
La longueur totale du contour est $39,7$ m.
\item Au centimètre, cela donne $39,70$ m.
\end{enumerate}
}

\corrige{27}{
\begin{enumerate}[label=\arabic*)]
\item Premier nombre :
\[
4,8\times 0,5=2,4.
\]
Deuxième nombre :
\[
\frac{3}{4}\times 60=3\times 15=45.
\]
Pour le troisième, on compare :
\[
7,09=7,090,\quad 7,9=7,900,\quad 7,099=7,099,\quad 7,901=7,901.
\]
Le plus petit est $7,09$.
Le code est donc :
\[
2,4\ ;\ 45\ ;\ 7,09.
\]
\item Les justifications sont les calculs précédents.
\item Ordre croissant :
\[
7,09<7,099<7,9<7,901.
\]
\end{enumerate}
}

\corrige{28}{
\begin{enumerate}[label=\arabic*)]
\item Total :
\[
240+150+90+120=600.
\]
Le total est cohérent.
\item Pourcentage bus :
\[
\frac{240}{600}=\frac{40}{100}=40\%.
\]
\item Pourcentage vélo :
\[
\frac{90}{600}=\frac{15}{100}=15\%.
\]
\item Un quart de $600$ vaut :
\[
600\div 4=150.
\]
La voiture représente donc un quart des élèves.
\item Tableau :
\[
\begin{array}{|c|c|c|}
\hline
\text{Moyen} & \text{Nombre} & \text{Pourcentage}\\
\hline
\text{Bus} & 240 & 40\%\\
\hline
\text{Voiture} & 150 & 25\%\\
\hline
\text{Vélo} & 90 & 15\%\\
\hline
\text{Marche} & 120 & 20\%\\
\hline
\end{array}
\]
La somme des pourcentages vaut $100\%$.
\end{enumerate}
}

\corrige{29}{
\begin{enumerate}[label=\arabic*)]
\item Le schéma est un rectangle séparé en deux rectangles identiques par un axe passant par les milieux des deux grands côtés.
\item Aire totale :
\[
A=10\times 6=60.
\]
L'aire totale est $60$ m$^2$.
\item Chaque moitié a pour aire :
\[
60\div 2=30.
\]
Chaque partie mesure $30$ m$^2$.
\item Les deux parties sont symétriques. La symétrie conserve les longueurs et les aires, donc les deux parties ont la même aire.
\item Dans une moitié :
\[
30\times 4=120.
\]
Il faut $120$ fleurs.
\end{enumerate}
}

\corrige{30}{
\begin{enumerate}[label=\arabic*)]
\item Distance en $1$ minute :
\[
7,5\div 3=2,5.
\]
Le robot parcourt $2,5$ m par minute.
\item En $48$ minutes :
\[
48\times 2,5=120.
\]
Il peut parcourir $120$ m.
\item La roche est à $125$ m. Comme $120<125$, il ne peut pas l'atteindre.
\item Distance manquante :
\[
125-120=5.
\]
Comme il parcourt $2,5$ m par minute :
\[
5\div 2,5=2.
\]
Il lui manque $2$ minutes.
\item Conclusion : avec une batterie pleine, le robot parcourt $120$ m en $48$ minutes. La roche étant située à $125$ m, il lui manque $5$ m, soit $2$ minutes de fonctionnement.
\end{enumerate}
}

\corrige{31}{
\begin{enumerate}[label=\arabic*)]
\item Périmètre complet :
\[
P=2\times(14+8)=2\times 22=44.
\]
Le périmètre complet est $44$ m.
\item On enlève une ouverture de $3$ m :
\[
44-3=41.
\]
Il faut $41$ m de barrière.
\item Prix de la barrière :
\[
41\times 6,20=254,20.
\]
\item Aire :
\[
A=14\times 8=112.
\]
L'aire est $112$ m$^2$.
\item Chaque dalle couvre $2$ m$^2$ :
\[
112\div 2=56.
\]
Il faut $56$ dalles.
\item Prix des dalles :
\[
56\times 9,50=532.
\]
\item Budget total :
\[
254,20+532=786,20.
\]
Le budget total est $786,20$ euros.
\end{enumerate}
}

\corrige{32}{
\begin{enumerate}[label=\arabic*)]
\item Total :
\[
8+6+4+2=20.
\]
Il y a $20$ cristaux.
\item Doré :
\[
P(D)=\frac{2}{20}=\frac{1}{10}=0,1=10\%.
\]
\item Bleu :
\[
P(B)=\frac{8}{20}=\frac{2}{5}=0,4=40\%.
\]
\item Ne pas tirer vert : il y a $20-4=16$ cristaux non verts :
\[
P(N)=\frac{16}{20}=\frac{4}{5}=0,8=80\%.
\]
\item Les écritures ont été données :
\[
D:\frac{1}{10}=0,1=10\%,
\quad
B:\frac{2}{5}=0,4=40\%,
\quad
N:\frac{4}{5}=0,8=80\%.
\]
\item Ordre du moins probable au plus probable :
\[
D<B<N.
\]
\end{enumerate}
}

\corrige{33}{
\begin{enumerate}[label=\arabic*)]
\item Prix des billets :
\[
28\times 12,50=350.
\]
Les billets coûtent $350$ euros.
\item Prix des goûters :
\[
28\times 2,40=67,20.
\]
\item Coût total :
\[
350+180+67,20=597,20.
\]
\item Coût moyen par élève :
\[
597,20\div 28=21,328\ldots
\]
Soit environ $21,33$ euros par élève.
\item Aire du bassin :
\[
A=25\times 12=300.
\]
Le bassin a une aire de $300$ m$^2$.
\item Périmètre de la zone circulaire :
\[
P=\pi\times D\approx 3,14\times 8=25,12.
\]
\item Oui, le responsable a raison : $25,12$ m est très proche de $25$ m.
\end{enumerate}
}

\corrige{34}{
\begin{enumerate}[label=\arabic*)]
\item Tir à l'arc :
\[
\frac{2}{5}\times 360=2\times(360\div 5)=2\times 72=144.
\]
Il y a $144$ visiteurs au tir à l'arc.
\item Château :
\[
30\%=\frac{30}{100}.
\]
Donc :
\[
\frac{30}{100}\times 360=108.
\]
Il y a $108$ visiteurs au château.
\item Marché :
\[
360-144-108=108.
\]
Il y a $108$ visiteurs au marché.
\item Dimensions réelles :
\[
6\times 4=24\text{ m},
\qquad
3,5\times 4=14\text{ m}.
\]
\item Aire réelle :
\[
A=24\times 14=336.
\]
La place a une aire de $336$ m$^2$.
\item Chaque stand occupe $12$ m$^2$ :
\[
336\div 12=28.
\]
On peut installer au maximum $28$ stands.
\item Comme $28\geq 25$, la place est suffisante pour installer $25$ stands.
\end{enumerate}
}

\newpage

%====================================================
\section{Fiche méthode finale}
%====================================================

\begin{methodebox}
\textbf{Méthode 1 -- Lire un nombre décimal.}
\begin{enumerate}[label=\arabic*)]
\item Identifier la partie entière.
\item Identifier la partie décimale.
\item Repérer les rangs : dixièmes, centièmes, millièmes.
\item Pour comparer, ajouter des zéros si nécessaire.
\end{enumerate}
\end{methodebox}

\begin{methodebox}
\textbf{Méthode 2 -- Calculer avec des fractions.}
\begin{enumerate}[label=\arabic*)]
\item Pour une fraction d'une quantité, diviser par le dénominateur puis multiplier par le numérateur.
\item Pour additionner deux fractions, chercher un dénominateur commun.
\item Pour comparer deux fractions, les mettre au même dénominateur si besoin.
\end{enumerate}
\end{methodebox}

\begin{methodebox}
\textbf{Méthode 3 -- Résoudre un problème de proportionnalité.}
\begin{enumerate}[label=\arabic*)]
\item Vérifier que la situation est proportionnelle.
\item Organiser les données dans un tableau.
\item Utiliser le retour à l'unité ou un coefficient multiplicateur.
\item Rédiger une phrase de conclusion avec l'unité.
\end{enumerate}
\end{methodebox}

\begin{methodebox}
\textbf{Méthode 4 -- Résoudre un problème de géométrie.}
\begin{enumerate}[label=\arabic*)]
\item Faire une figure à main levée.
\item Coder les informations connues.
\item Choisir la bonne propriété : somme des angles, médiatrice, symétrie, périmètre, aire.
\item Écrire les calculs proprement.
\item Conclure par une phrase.
\end{enumerate}
\end{methodebox}

\begin{methodebox}
\textbf{Méthode 5 -- Calculer une probabilité simple.}
\begin{enumerate}[label=\arabic*)]
\item Compter le nombre total de cas possibles.
\item Compter le nombre de cas favorables.
\item Écrire :
\[
P=\frac{\text{cas favorables}}{\text{cas possibles}}.
\]
\item Simplifier la fraction si possible.
\item Convertir en décimal ou en pourcentage si demandé.
\end{enumerate}
\end{methodebox}

\begin{erreurbox}
\textbf{Erreurs classiques à éviter.}
\begin{itemize}
\item Oublier l'unité dans une réponse.
\item Confondre périmètre et aire.
\item Comparer $7,9$ et $7,09$ sans écrire $7,90$.
\item Écrire une probabilité supérieure à $1$.
\item Acheter un nombre décimal d'objets indivisibles : on achète $4$ pots, pas $3,9$ pots.
\item Additionner des fractions sans mettre au même dénominateur.
\item Croire que toutes les situations avec deux grandeurs sont proportionnelles.
\end{itemize}
\end{erreurbox}

\begin{retenirbox}
\textbf{Formules importantes.}
\[
P_{\text{rectangle}}=2\times(L+\ell)
\]
\[
A_{\text{rectangle}}=L\times \ell
\]
\[
A_{\text{carré}}=c\times c
\]
\[
P_{\text{cercle}}=\pi\times D=2\times \pi\times R
\]
\[
a\%=\frac{a}{100}
\]
\[
P(\text{événement})=\frac{\text{nombre de cas favorables}}{\text{nombre de cas possibles}}
\]
\[
\widehat{A}+\widehat{B}+\widehat{C}=180^\circ
\]
\end{retenirbox}

\begin{vigilancebox}
\textbf{Stratégie générale en contrôle.}
\begin{enumerate}[label=\arabic*)]
\item Lire l'énoncé deux fois.
\item Souligner les données utiles.
\item Identifier le thème : fraction, aire, proportionnalité, probabilité, géométrie.
\item Faire un schéma ou un tableau si nécessaire.
\item Écrire les calculs avec les unités.
\item Vérifier l'ordre de grandeur.
\item Répondre par une phrase.
\end{enumerate}
\end{vigilancebox}

\end{document}